%============================================================
% BAB 4: EKSPERIMEN DAN ANALISIS
% Kerangka siap isi (tanpa referensi gambar dummy yang dihapus).
%============================================================
\chapter{EKSPERIMEN DAN ANALISIS}

Uraian pada bab ini meliputi parameter eksperimen, karakteristik data, tempat dan waktu uji coba, spesifikasi peralatan, hasil pengujian sistem (unit test, kontrak API, isolasi multi-tenant, integrasi Midtrans, serta evaluasi MCP dan antarmuka Flutter Web), dan analisis ketercapaian tujuan Proyek Akhir.

\section{PARAMETER EKSPERIMEN}

Uraian parameter pengujian perangkat lunak dan skenario evaluasi modularitas serta isolasi data antar-gedung akan diisi pada tahap implementasi.

\section{KARAKTERISTIK DATA}

Uraian karakteristik data uji (data gedung, kamar, kontrak huni bulanan, tagihan, tiket kerusakan, dan inventaris) akan diisi pada tahap implementasi.

\section{TEMPAT DAN WAKTU UJICOBA}

Uraian lokasi, lingkungan pengujian (Docker Compose, PostgreSQL, Keycloak OIDC, NestJS, dan Flutter Web), serta jadwal pelaksanaan uji coba akan diisi pada tahap implementasi.

\section{SPESIFIKASI PERALATAN UJICOBA}

Uraian spesifikasi perangkat keras dan perangkat lunak yang digunakan selama pengembangan dan pengujian sistem akan diisi pada tahap implementasi.

\section{HASIL EKSPERIMEN}

Hasil pengujian fungsional, pengujian kombinasi modul, pengujian kontrak REST API, pengujian isolasi tenant, serta evaluasi pengguna akan disajikan pada bagian ini.

\section{ANALISIS HASIL EKSPERIMEN}

Analisis hasil pengujian dan pemetaannya terhadap rumusan masalah serta tujuan pada Bab~I akan diuraikan pada bagian ini.
